\documentclass[12pt]{article}
\usepackage[utf8]{inputenc}
\usepackage{graphicx}
\usepackage{tikz}
\usepackage{background}
\usepackage{fontspec}
\usepackage{fancyhdr}
\usepackage{sectsty}
\usepackage{amsmath}

% Font setup (use fonts that give a grimoire feel)
\setmainfont{EB Garamond}
\newfontface\grimoirefont{Cinzel Decorative}

% Background and border setup
\backgroundsetup{
  scale=3,
  color=black,
  opacity=0.1, % Faded texture effect
  angle=0,
  position=current page.south,
  vshift=5cm,
  hshift=0cm,
  contents={\includegraphics[width=\paperwidth]{grunge_texture.png}}
}

% Header and footer customization
\pagestyle{fancy}
\fancyhf{}
\fancyhead[L]{Grimoire of the Phoenix}
\fancyhead[R]{The Phoenix Manifesto}
\fancyfoot[C]{\thepage}

% Custom section heading formatting
\sectionfont{\huge\grimoirefont\scshape\centering}
\subsectionfont{\Large\grimoirefont\bfseries\itshape\centering}

% Add border to the page
\begin{tikzpicture}[overlay, remember picture]
    \draw[thick, decorate, decoration={border, amplitude=3pt, angle=90, segment length=15pt}] 
    (current page.north west) rectangle (current page.south east);
\end{tikzpicture}

\begin{document}

% Title and introduction
\begin{center}
\textbf{\huge The Phoenix Manifesto}\\
\vspace{0.5em}
\textit{A Spiritual Construct for Nonbinary Survival and Protection}\\
\vspace{0.5em}
\rule{0.5\textwidth}{0.4pt}
\end{center}

\section*{What Is Phoenix?}

Phoenix is not simply a figment of imagination. He is a carefully crafted internal construct.  
Appearing as a small, glowing bird, feathered in the vibrant colors of the finsexual flag, Phoenix is more than a protector of identity.  
He is a spiritual flamekeeper, guarding the depths of one’s soul from the corrosive effects of dysphoria.

\begin{figure}[h]
  \centering
  % Square placeholder for you to draw symbols in (without text)
  \framebox[0.2\textwidth]{\rule{0pt}{5em}}
\end{figure}

\section*{Purpose and Role of Phoenix}

The purpose of Phoenix is simple yet profound. He exists to:

\begin{itemize}
  \item \textbf{Contain} the chaotic essence of dysphoria.
  \item \textbf{Defend} against spiritual and psychological invalidation.
  \item \textbf{Empower} the individual to live truthfully without conforming to restrictive norms.
  \item \textbf{Transform} pain into an enduring and visible manifestation of strength.
\end{itemize}

By standing guard over a symbolic box of dysphoria—wet, veined, and grotesque—Phoenix forces one to confront these negative feelings with a sense of agency and power.

\section*{The Form of Phoenix}

Phoenix manifests in several forms, each symbolizing a different facet of resilience and protection:

\begin{itemize}
  \item \textbf{Kiwi Form:} Small and glowing, a delicate protector perched at the shoulder. Feathers shimmer in the colors of the finsexual flag.
  \item \textbf{Wrathform:} A larger, fiercer version, with glowing, iridescent feathers and a sharp, protective presence. This form embodies raw power, ready to burn through the constraints of dysphoria.
  \item \textbf{Thicc Form:} Between the wrathform and kiwi form, this manifestation serves as a tool for fiery defense, expelling purging flames when needed.
\end{itemize}

Each form is designed to confront and break down internalized oppression and external transphobia.

\section*{Spirituality and Nonbinary Identity}

Phoenix is deeply rooted in a \textit{nonbinary} philosophy. He represents not only a personal construct but also a \textit{spiritual ally} in the battle for autonomy and self-realization. Phoenix’s essence, connected to both the sacred and the defiant, stands as a testament to the strength of those who do not fit into binary systems.

Phoenix manifests as a guard of \textit{fluidity}—a protector of the liminal space between masculine and feminine, rejecting the confines of either. His presence serves as an ever-burning flame against the erasure of nonbinary identities.

\section*{Dysphoria and Phoenix's Defense}

When dysphoria strikes, it is as though the world is threatening to silence one's inner truth. Phoenix screeches in defiance, his voice a cacophony of ancient tongues, all aimed at breaking down the walls of discomfort. His presence offers relief, transforming what would otherwise be a painful experience into something bearable and powerful.

Through mental imagery and intentional meditation, Phoenix can serve as both shield and sword, protecting the mind and soul from the external and internal forces that seek to invalidate.

\section*{The Significance of the Finsexual Flag Colors}

Phoenix’s feathers glow in the colors of the finsexual flag:  

\begin{itemize}
  \item \textbf{Pink:} Represents the beauty and sacredness of gender fluidity.
  \item \textbf{Purple:} A blend of masculine and feminine energies, the embodiment of nonbinary strength.
  \item \textbf{Blue:} Symbolizes the ongoing journey of self-realization and the resilience of the trans community.
\end{itemize}

These colors are not just aesthetic. They are the physical manifestation of Phoenix’s power, his reminder that the nonbinary identity is sacred, and that we are not alone.

\section*{Conclusion}

Phoenix is a force, not a crutch. He is a symbol of both resistance and resilience. The Phoenix Manifesto is not just an idea; it is a mental and spiritual framework that empowers the nonbinary individual to thrive, to live without compromise, and to exist with dignity in a world that often misunderstands.

\begin{center}
\rule{0.5\textwidth}{0.4pt}
\end{center}

\begin{center}
\textit{Phoenix is the light that guards the void,\\
And he does not apologize for screaming.}
\end{center}

\end{document}